\vspace{0.5cm}
\section{Introducción}
\justify
La planificación financiera de proyectos de construcción en Paraguay depende, en parte importante, de poder estimar con anticipación el costo de los materiales. Esa estimación resulta difícil porque los precios del cemento y el ladrillo no se mueven de manera predecible: responden a una combinación de factores que incluye desde las condiciones logísticas del transporte fluvial hasta los efectos de eventos sanitarios de alcance global. La consecuencia práctica es que muchas empresas del sector trabajan con márgenes de incertidumbre que complican la viabilidad de los proyectos.

\justify
Este contexto se volvió especialmente evidente en dos episodios recientes. El primero fue la pandemia de COVID-19: las restricciones de circulación y el parate en la cadena de suministros generaron incrementos notorios en los precios de materiales esenciales, tal como documentó la revista Unida Científica (citada en Zuba, 2021). El segundo fue la bajante histórica del río Paraguay en 2021, que forzó a la Industria Nacional del Cemento a sustituir el transporte fluvial por terrestre para mantener el abastecimiento de materia prima, decisión que tuvo costos directos sobre el precio al consumidor final.

\justify
Ambos episodios evidencian que los modelos descriptivos tienen un alcance limitado para orientar decisiones prospectivas. Se requieren enfoques que integren la heterogeneidad de los factores que inciden en los precios y que sean capaces de proyectar escenarios futuros con un nivel de detalle útil para la toma de decisiones. Este trabajo propone una alternativa concreta: combinar arquitecturas de aprendizaje profundo, LSTM y GRU, con el modelo estadístico clásico SARIMAX, usando datos públicos y herramientas de software libre, para predecir la evolución mensual de los precios del cemento y el ladrillo común en Paraguay.

\section{Justificación y Motivación}
\justify
La literatura sobre precios de materiales de construcción en Paraguay y América Latina muestra un patrón llamativo: abundan los análisis descriptivos de lo que ya ocurrió, pero escasean los modelos que permitan anticipar lo que viene. Esa brecha no es menor. Para una empresa constructora, conocer con razonable anticipación que el precio del cemento podría subir en los próximos meses puede definir si un proyecto se adjudica a un precio sostenible o si termina generando pérdidas.

\justify
Otro vacío identificado es el tratamiento de variables externas. El nivel del río Paraguay y los períodos de confinamiento son factores que claramente afectan la logística y la disponibilidad de insumos, pero raramente se incorporan de forma explícita en los modelos predictivos existentes. Incluirlos no es solo un refinamiento técnico: es reconocer que los precios responden a la realidad geográfica y socioeconómica del país.

\justify
Este trabajo parte de esa doble motivación. Por un lado, busca generar evidencia concreta sobre qué tan bien pueden predecirse los precios de estos materiales cuando se utilizan modelos de aprendizaje profundo frente a enfoques estadísticos clásicos. Por otro, propone una metodología reproducible —construida sobre datos públicos y herramientas de acceso libre— que pueda ser adoptada o extendida tanto por investigadores como por actores del sector privado y público.

\section{Objetivos}
\subsection{Objetivo General}
\justify
Desarrollar y comparar modelos predictivos basados en métodos estadísticos y de aprendizaje profundo para estimar la evolución de los precios del cemento y el ladrillo común en Paraguay, incorporando variables exógenas como el confinamiento por COVID-19 y el nivel del río Paraguay.

\subsection{Objetivos Específicos}
\begin{itemize}
\item Recopilar y procesar series temporales de precios de materiales de construcción y del nivel del río Paraguay para el período enero 2014 – julio 2025.
\item Implementar un modelo LSTM univariado para la predicción del nivel mensual del río Paraguay.
\item Desarrollar modelos SARIMAX, LSTM y GRU multivariados para la predicción de precios de materiales de construcción.
\item Evaluar el desempeño de los modelos mediante métricas de error (RMSE) y comparar su capacidad predictiva bajo distintos escenarios experimentales (con y sin confinamiento por COVID-19).
\item Analizar el impacto de las variables exógenas sobre la dinámica de precios y discutir su relevancia en el contexto económico y logístico nacional.
\end{itemize}

\section{Organización del Libro}
\begin{itemize}
\item \textbf{Capítulo 2:} presenta los trabajos relacionados, incluyendo una revisión de investigaciones previas sobre la predicción de precios de materiales de construcción y las limitaciones identificadas en dichos estudios.
\item \textbf{Capítulo 3:} desarrolla el fundamento teórico de los modelos empleados, describiendo los conceptos matemáticos de ARIMA, SARIMA, SARIMAX y las arquitecturas de redes neuronales recurrentes LSTM y GRU.
\item \textbf{Capítulo 4:} describe la metodología aplicada, detallando la recolección y preprocesamiento de datos, la formulación del modelo LSTM univariado para el nivel del río Paraguay y los modelos predictivos de precios (SARIMAX, LSTM y GRU).
\item \textbf{Capítulo 5:} expone los resultados obtenidos, tanto en la predicción del nivel del río como en la estimación de precios de los dos materiales analizados.
\item \textbf{Capítulo 6:} desarrolla la discusión comparativa de los resultados, evaluando el desempeño de cada modelo y su capacidad predictiva frente a las variables exógenas consideradas.
\item \textbf{Capítulo 7:} presenta las conclusiones generales y propone líneas de trabajo futuro orientadas a la mejora del enfoque predictivo y su extensión a otros materiales o indicadores económicos del sector construcción.
\end{itemize}
